% IV. Програмна реализация.
% Пише се докато кодът се пише, не преди това. Твърденията тук трябва да сочат
% към съществуващ файл.
\section{Програмна реализация}
\label{sec:implementation}

\subsection{Обща структура на изходния текст}

Реализацията следва слоевете от Раздел~\ref{sec:design}. Всеки слой е отделна
преводна единица с публичен заглавен файл, а зависимостите между тях са
еднопосочни и се проверяват от системата за построяване. Външните зависимости са
две: пакетът за оптимизация и библиотеката за четене и записване на JSON.

\TODO{таблица на модулите: име, отговорност, публичен интерфейс; попълва се от
действителното дърво на изходния текст, не от плана}

\subsection{Слой на модела}

Слоят на модела съдържа представянето на графа и проверката на инвариантите му.
Проверката за ацикличност се извършва с топологично подреждане при зареждане на
сценария. Същото подреждане се преизползва при изчисляването на най-ранните
срокове, така че обходът на графа се прави веднъж, а не два пъти.

\TODO{коментар на изходния текст на топологичното подреждане и на изчисляването
на сроковете, с листинг}

\subsection{Слой на оценяването}

Двата метода за многокритериална оценка са зад общ интерфейс, който приема
матрицата проекти по критерии и връща вектор от оценки. Аналитичната йерархия
изчислява теглата от матрицата на попарните сравнения и отчита коефициента на
съгласуваност; при стойност над приетия праг оценяването се отказва вместо да
върне число, изведено от несъгласувани преценки. TOPSIS нормализира матрицата,
изчислява разстоянията до идеалната и до най-лошата точка и връща относителната
близост.

\TODO{запиши приетия праг за коефициента на съгласуваност и източника, от който е
взет; не го попълвай по памет}

\subsection{Слой на подбора}

Задачата се формулира с двоични променливи, по една на проект, целева функция
сумата от оценките и ограничения по бюджет, по капацитет за всеки период и по
предхождане. Лакомата евристика е реализирана в същия слой и се изпълнява
винаги, за да служи за долна граница при проверката на резултата от решателя.

\subsection{Слой на представянето}

Слоят на представянето приема сценарий, изпълнява останалите три слоя и връща
резултата. Обяснението към препоръката се сглобява тук: за всеки отхвърлен
проект се записва кое ограничение е било определящо.

\TODO{опиши крайните точки на интерфейса и формата на отговорите, след като са
реализирани}

\subsection{Тестове}

Проверката на реализацията е разделена на два вида. Модулните тестове покриват
отделните изчисления: топологичното подреждане, нормализацията в TOPSIS,
изчисляването на теглата. Тестовете върху цялата система изпълняват сценарии с
известно предварително решение и сравняват изхода с него.

\TODO{посочи достигнатото покритие и инструмента, с който е измерено}
